\thispagestyle{empty}
\vspace*{2cm}

\begin{center}
\textbf{\LARGE The BX3 Framework: Architectural Accountability for Recursive Multi-Agent Systems}\\[1.5em]

\textbf{Jeremy Beebe}$^{1}$\footnote{Corresponding author. Email: jeremy@bxthre3.com} \\[0.6em]
\textbf{$^{1}$Bxthre3 Inc, Monte Vista, Colorado, USA}\\[2em]

\textbf{Abstract}\\[0.5em]
\end{center}

\noindent
We present the \textbf{BX3 Framework} (Behavior, Execution, and eXchange Architecture) ---
a novel architectural framework for recursive multi-agent AI systems that enforces
human accountability as a compulsory structural property, not a policy aspiration.
The framework organizes every node into three immutable functional layers ---
\textbf{Purpose} (intent and final authority), \textbf{Bounds} (bounded reasoning within
provisioned operating ranges), and \textbf{Fact} (deterministic enforcement and forensic
auditability) --- and introduces five named architectural pillars that collectively
constitute an \textbf{upstream accountability guarantee}: that BX3 systems
fail upward into human consciousness, never downward into autonomous chaos.

The five pillars are: (1) \textbf{Loop Isolation}, which prevents state circularity
in recursive execution graphs; (2) \textbf{Recursive Spawning}, which maintains immutable
parent-child accountability chains as nodes spawn subordinate agents; (3)
\textbf{Spatial Firewall}, which enforces execution domain separation between spawned
sub-agents; (4) \textbf{Sandbox Gate}, which provides controlled inter-domain
communication channels with mandatory authorization verification; and
(5) \textbf{Bailout Protocol}, which compels every unresolved exception state to
propagate upward to a named human principal, bypassing all intermediate machine
actors, with immutable forensic attribution at every step.

We provide formal specifications for each pillar, establish the correctness
properties they jointly guarantee, and describe the complete BX3 three-layer
implementation model.  Evaluated against the Tiered Agentic Oversight (TAO) model,
the BX3 Framework provides stronger accountability guarantees: TAO permits
absorption at intermediate machine tiers; the BX3 Framework forbids it
architecturally.  We further characterize the framework's properties under
adversarial conditions and identify the directed research agenda required to
extend these guarantees to high-stakes production deployments.

\noindent\textbf{Keywords:} multi-agent systems, accountability, human-in-the-loop,
fail-safe design, recursive AI, architectural frameworks, BX3, Bailout Protocol

\newpage
\tableofcontents
\newpage